We investigate automatic, joint style and data optimization for MapLibre-based map rendering pipelines, asking whether composing style-level and data-level transformations yields additive or synergistic gains.
Interactive web maps built on vector tiles face a complex performance landscape: every frame requires fetching tiles over the network, decoding them, evaluating style expressions for each feature, and submitting draw calls to the graphics \ac{API}.
Despite this, style authoring and tile data preparation have traditionally been treated as independent concerns, leaving the interaction between the two unexploited.

We design and implement a standalone optimizer that operates at three levels.
Expression rewrites come first: a set of peephole rules - constant folding, algebraic simplification, stats-driven folding from tile statistics, and selectivity-based predicate reordering - is applied to the style's \ac{JSON} expressions via fixpoint iteration.
On top of that, structural passes work on a typed representation of the full style document to do dead layer elimination, metadata refinement, layer merging, and cleanup.
The third level leaves the style behind and reaches into the tile data: a pruning advisory derived from the optimized style removes unused properties, geometry types, and feature values, and the surviving content is re-encoded into the columnar MapLibre Tiles format with automatic encoding strategy selection.

We evaluated across 15~styles, 18~scenarios, and 19~cumulative ablation steps plus one plain-tile-shaving comparison, and we observe that tile shaving dominates load-time improvements, achieving load-time reductions of \qtyrange{63}{74}{\percent} (median \qty{71}{\percent}) and \ac{FPS} gains of \qtyrange{100}{200}{\percent} (median \qty{152}{\percent}) across the ten deep-corpus styles we benchmarked end-to-end; across the 11 \ac{OMT}-compatible styles for which we measured tile shaving without browser rendering, shaving effectiveness varied \qtyrange{19.8}{46.3}{\percent}.
Our new \ac{MLT} encoder reduces total tile volume by \qty{28}{\percent} over the reference encoder and \qty{12}{\percent} over gzip-compressed \ac{MVT} on the Germany OpenMapTiles tileset.
Style-level optimizations alone reduce style transfer size, primarily benefiting initial load: across 13 unseen styles in a static-only generalization corpus the static passes produce a median \qty{31.4}{\percent} raw reduction (after gzip median \qty{11.4}{\percent}, range \qtyrange{3.9}{22.4}{\percent}) with identical rendering output.
We find that the two levels compose additively rather than synergistically at the aggregate level; per-scenario synergy appears in 17 of 36 style-scenario pairs but is too sparse to lift the aggregate above the additive expectation.
